\section{Book Architecture}

The manuscript is organized around nine teaching jobs:
\begin{enumerate}[nosep]
  \item Chapters~\ref{ch:formal-language} and \ref{ch:dynamics} define the
        language of the model.
  \item Chapter~\ref{ch:planning} explains why planning is structurally
        necessary.
  \item Chapter~\ref{ch:consciousness-biology} introduces the empirical
        consciousness layer and ties state transitions to biological evidence.
  \item Chapter~\ref{ch:emergence} explains how durable frameworks arise from
        repeated interaction, not from a single isolated choice.
  \item Chapter~\ref{ch:advanced-dynamics} upgrades the mathematical language
        from a lightweight barrier model to a richer state-space view.
  \item Chapter~\ref{ch:cases} makes the model concrete across multiple
        domains.
  \item Chapter~\ref{ch:interventions} turns analysis into design patterns.
  \item Chapter~\ref{ch:exercises} tests whether the reader can reuse the
        framework.
  \item The appendices preserve notation, templates, and clearly isolated
        frontier hypotheses.
\end{enumerate}
